\paragraph*{In short}

\begin{table}[htbp]
\centering\small
\begin{tabularx}{\linewidth}{@{}lXX@{}}
\toprule
\textbf{Syntax} & \textbf{What it does} & \textbf{Runs on} \\
\midrule
\texttt{\_\allowbreak{}\_\allowbreak{}device\_\allowbreak{}\_\allowbreak{} fn()} & Device function; passed through unchanged to the generated source & Device \\
\texttt{call fn(args)} & Invoke a \texttt{\_\_device\_\_} function from a \texttt{\_\_co\_\_} body & Device \\
\texttt{\_\allowbreak{}\_\allowbreak{}cpp\_\allowbreak{}\_\allowbreak{}(".\allowbreak{}.\allowbreak{}.\allowbreak{}")} & Inject verbatim C++ at the splice point (ordinary string) & Target \\
\texttt{\_\allowbreak{}\_\allowbreak{}cpp\_\allowbreak{}\_\allowbreak{}(R"(.\allowbreak{}.\allowbreak{}.\allowbreak{})")} & Inject verbatim C++ (raw string; use for \texttt{asm volatile}) & Target \\
\texttt{\#\allowbreak{}define NAME value} & Object-like macro; shared across \texttt{\_\_co\_\_} and host regions & Preprocessor \\
\texttt{\#if} / \texttt{\#ifdef} / \texttt{\#error} & Conditional compilation and compile-time assertions & Preprocessor \\
\texttt{-DNAME=value} & Define macros from the command line for tile sweeps & Preprocessor \\
\bottomrule
\end{tabularx}
\end{table}

Section~\ref{sec:3-11} turns to the other side of the workflow: what to do when the kernel compiles but the output is wrong --- debugging, verbose modes and systematic narrowing. The two worked programs in Section~\ref{sec:3-10} put the interoperability mechanisms into context.
